\documentclass[acmsmall,11pt,screen,nonacm]{acmart}
\setcopyright{none}
\settopmatter{printacmref=false,printccs=false,printfolios=true}
\geometry{letterpaper,margin=1in,headheight=13pt,footskip=24pt,
          includehead=false,includefoot=false,heightrounded=false}
\usepackage{mathtools}
\renewcommand{\bibfont}{\normalfont\normalsize}
\newtheorem{theorem}{Theorem}
\newtheorem{lemma}[theorem]{Lemma}
\newtheorem{proposition}[theorem]{Proposition}
\newcommand{\OPT}{\operatorname{OPT}}
\newcommand{\E}{\mathbb{E}}
\newcommand{\Prb}{\mathbb{P}}
\newcommand{\bits}{\{0,1\}}
\title{A Near-Logarithmic Approximation for All-or-Nothing Flow}
\author{Zijian Zeng}
\affiliation{\institution{UCSI University}\city{Kuala Lumpur}\country{Malaysia}}
\email{1002266693@ucsiuniversity.edu.my}
\renewcommand{\shortauthors}{Zijian Zeng}
\begin{document}
\begin{abstract}
\normalsize
The all-or-nothing flow problem asks for a largest subset of terminal pairs that
can each send one splittable unit of flow without violating edge capacities.
For an undirected unit-capacity graph with $k$ indexed pairs, we give a randomized
polynomial-time approximation with factor
$C\log(k+2)\exp(C\sqrt{\log\log(k+4)})$ for an absolute constant $C$.
The output is always feasible, and the cardinality guarantee holds with
probability at least $3/4$. This gives the $\log^{1+o(1)}k$ guarantee asked for
in Conjecture~5.1 of Bansal, Jambulapati, and Saranurak.

The reduction combines their integral-node-weight flow-expander decomposition
with the established congestion-one matching rounding of Chekuri, Khanna, and
Shepherd. A fixed family of uniformly weighted path copies permits a single
decomposition. Original unit terminal edges are retained as bottlenecks, so a
feasible selection of copies admits distinct original pairs. An exact
small-parameter branch makes the final guarantee depend on $k$, even when the
graph is much larger than the terminal set. Sampling uses a strict polynomial
bound on random bits.
\end{abstract}
\maketitle
\pagestyle{plain}
\raggedbottom
\clearpage

\section{The result and its context}

Let $G=(V,E)$ be a finite undirected graph with unit edge capacities, and let
$(s_i,t_i)$, $i\in[k]$, be indexed terminal pairs. A feasible solution chooses
$I\subseteq[k]$ and routes one unit between $s_i$ and $t_i$ for each $i\in I$.
Each commodity may split among paths, but the sum of all path flows on each
edge is at most one. The objective is $|I|$; denote its optimum by $\OPT$.
Shared endpoints, disconnected inputs, and indexed duplicate pairs are allowed.
A pair with coincident endpoints has a zero-edge route.

Chekuri, Khanna, and Shepherd give an $O(\log^2 k)$ approximation for the
edge-capacitated problem~\cite[Theorem~3.2]{CKS05}. Bansal, Jambulapati, and
Saranurak obtain flow-expander decomposition with nearly logarithmic overhead,
and explicitly ask whether a $\log^{1+o(1)}k$ approximation for all-or-nothing
flow is possible~\cite[Conjecture~5.1]{BJS26}. Their discussion explains why
changing between throughput and concurrent-flow subproblems prevents an
immediate substitution into the earlier framework. The theorem below supplies
an admission-preserving interface for a fixed flow instead.

\begin{theorem}\label{thm:main}
There is one randomized algorithm, polynomial in the explicit input bit
length, for all the inputs above. It always returns a feasible subset and
unit flows of congestion at most one. For an absolute constant $C$, with
probability at least $3/4$ its cardinality is at least
\[
 \frac{\OPT}{C\log(k+2)\exp\!\bigl(C\sqrt{\log\log(k+4)}\bigr)}.
\]
Its number of random bits is polynomially bounded. Repetition and selection
of the largest feasible output amplify the success probability.
\end{theorem}

Here and below logarithms in asymptotic bounds are natural; explicitly written
$\log_2$ is binary. The denominator is $\log^{1+o(1)}k$ as $k\to\infty$.
Every admitted pair receives a full unit, and there is no congestion allowance.

\paragraph{What the reduction contributes.}
Uniformization and terminal duplication already occur in
\cite[Lemma~1.1]{CKS05}, and the decomposition and matching rounding used here
are prior results. The reduction keeps the original unit edges for terminal
occurrences while adding many copy terminals beyond them. These original
edges prevent two copies of the same pair from receiving integral admission.
A fixed feasible path family then charges decomposition losses to edges only
once. The copy scale cancels against the rounding yield. Finally, a small-$k$
exact branch bounds the number of copy terminals by $O(k^2)$ in the other
branch. These steps yield the full general-input improvement, rather than
an improvement for an auxiliary routing problem.

Anand subsequently gives improved flow-expander overhead and node-weighted
extensions~\cite[Theorems~1.4 and~5.1]{Anand26}. We use the precise
integral-total-weight guarantee of~\cite{BJS26}; we do not infer a stronger
$k$-dependent admission theorem from an $n$-dependent decomposition bound.
The algorithm below uses these established primitives to obtain the stated
admission guarantee. Codex assisted materially with the proof development,
source comparisons, and exposition; the final section describes these roles.

\section{Two established algorithmic ingredients}\label{sec:ingredients}

For a nonnegative node weighting $A$, write $|A|=\sum_v A(v)$.
An unordered-pair demand $D$ \emph{respects} $A$ if
$\sum_{v\ne u}D(\{u,v\})\le A(u)$ for every $u$. The weighting is
$\phi$-flow-expanding if every such demand can be routed with congestion
at most $1/\phi$.

Fix a nondecreasing, computable overhead bound
\begin{equation}\label{eq:gamma}
 \Gamma(a)=O\!\left(\log(a+2)
       \exp\!\bigl(O(\sqrt{\log\log(a+4)})\bigr)\right),
 \qquad \Gamma(a)\ge1,
\end{equation}
with constants sufficiently large for the following theorem. Values of
$\Gamma$ may be taken to be integers of polynomial bit length.

\begin{proposition}[Bansal--Jambulapati--Saranurak~\cite{BJS26}]
\label{prop:bjs}
Given an undirected unit-capacity graph, an integral node weighting $A\ge0$,
and $\phi>0$, a polynomial-time algorithm computes $C\subseteq E$ with
$|C|\le\phi|A|\Gamma(|A|)$. For each connected component $U$ of $G-C$,
the restriction $A|_U$ is $\phi$-flow-expanding in the induced graph $G[U]$.
\end{proposition}

This is Theorem~3.1 of~\cite{BJS26}; zero weights are permitted. If needed,
apply it separately to connected components with the same global $\phi$.
Monotonicity of $\Gamma$ bounds the sum of cut costs by the displayed global
bound. Components with total weight zero require no cut.

For $q=|X|$ distinct terminals, call $X$ \emph{$\alpha$-edge-flow-linked} if
the demand $\alpha/q$ on each unordered pair of distinct terminals is
routable with congestion one. This is the uniform specialization of the
weighted product-demand definition in~\cite{CKS05}.

\begin{proposition}[Chekuri--Khanna--Shepherd~\cite{CKS05,CKS13}]
\label{prop:cks}
There is an absolute $c_0>0$ such that, for $0<\alpha\le1$, a matching $M$
whose endpoints $X$ are $\alpha$-edge-flow-linked has a polynomial-time
congestion-one unit routing for at least $c_0\alpha|M|$ pairs of $M$.
\end{proposition}

The general matching statement is Theorem~3.1 of~\cite{CKS05}. Its footnote
requires strengthened clustering. Section~\ref{sec:correction} explains the
support from the corrected journal clustering in~\cite[Section~3.4]{CKS13}
for the leaf terminals used here. We do not apply the journal's nice-instance
Theorem~3.8 directly to arbitrary components. The conclusion is a splittable
unit routing, without an integral-path requirement.

\section{A fixed feasible family of uniform path copies}\label{sec:copies}

Admit all coincident-endpoint pairs and remove them for now. If no remaining
pair has connected endpoints, return these free admissions. In the rest of
the proof $k$ counts the remaining pairs; the bound for the initial $k$ is
no smaller after increasing its absolute constant.

For every terminal \emph{occurrence}, add a distinct leaf joined to its
original vertex by a unit edge, and move that terminal occurrence to the
leaf. Call these the \emph{original occurrence leaves and edges}. Even when
two indexed pairs have identical endpoints, their new leaves are distinct.
This normalization preserves $\OPT$: extend each admitted unit flow along
its own two new edges, or project a normalized flow after removing its
pendant segments. Simple-path decompositions suffice after cycle removal.

Choose an integer $N\ge16$ at least the encoding length and the numbers of
vertices, edges, and pairs of this normalized input. It can be chosen
polynomially bounded in the original explicit encoding length. Set
$b=\lceil\log_2N\rceil$. If $k\le b$, enumerate all subsets and test unit-flow
feasibility by a compact rational multicommodity-flow linear program.
There are $2^k\le2N$ subsets, so this branch is polynomial and returns an
optimum solution. Henceforth assume $k>b$.

Solve the fractional admission relaxation
\[
 \max F=\sum_i x_i,\qquad 0\le x_i\le1,
 \qquad \text{simultaneously route $x_i$ units for pair $i$}.
\]
Its compact edge formulation has polynomial size. Fix an optimum rational
solution and a polynomial-size decomposition into simple paths, cancelling
cycles first. Since some pair is connected,
\begin{equation}\label{eq:F}
 1\le F\le k,\qquad F\ge\OPT.
\end{equation}

Set
\begin{equation}\label{eq:scale}
 T=64\lceil\log_2(2N)\rceil,\qquad \varepsilon=\frac1{4T}.
\end{equation}
For commodity $i$, partition its path mass into
$d_i=\lceil Tx_i\rceil$ positive pieces, each of mass at most $1/T$;
take no pieces if $x_i=0$. Process the path list in order and split a path
when a piece becomes full. If $D=\sum_i d_i$, then
\begin{equation}\label{eq:pieces}
 D\le TF+k\le k(T+1).
\end{equation}
The number of positive fragments is at most the original number of paths
plus $D$. In particular, small positive masses do not require a number of
copies proportional to their denominators.

Independently in each piece, select one path fragment $P$ with probability
$T$ times its mass, or select nothing with the remaining probability.
For the moment this is the ideal rational sampling law. Let $K$ count all
selected paths and let $Z_e$ count those that use normalized edge $e$.
Independence across pieces and feasibility of the initial flow give
\begin{equation}\label{eq:means}
 \E K=TF,\qquad \E Z_e\le T.
\end{equation}
Each variable is a sum of independent Bernoulli variables, even when several
pieces belong to one original pair. Chernoff bounds imply
\begin{equation}\label{eq:chernoff}
 \Prb(K<TF/2)\le e^{-TF/8}\le e^{-T/8},\qquad
 \Prb(Z_e>4T)\le(e^3/4^4)^T<8^{-T}.
\end{equation}
For the second estimate, exponential Markov with parameter $\log4$ and
mean at most $T$ bounds even the event $Z_e\ge4T$. Since there are at most
$N$ normalized edges, $N\ge16$, and $T\ge320$,
$e^{-T/8}+N8^{-T}<1/32$.

\paragraph{A bounded sampler.}
For each piece list its rational outcomes, including a null outcome. Let
$Q$ be the sum of these list lengths and put
$B=\lceil\log_2(64Q)\rceil$. Draw $B$ unbiased bits per piece, yielding
$U\in\{0,\ldots,2^B-1\}$, and locate $U/2^B$ in its cumulative probability
intervals. The probability of an interval differs from its ideal length
by at most $2^{-B}$. Thus the total variation distance of the implemented
joint law from the ideal independent law is at most
$Q2^{-B}\le1/64$, by summing marginal total variation distances.
All comparisons are rational. This sampler has strict polynomial time and
random-bit bounds, with no rejection loop.

It follows that the event
\begin{equation}\label{eq:good}
 \mathcal G=\{K\ge TF/2\text{ and }Z_e\le4T\text{ for every normalized edge }e\}
\end{equation}
has probability greater than $1-1/32-1/64=61/64$ under the implemented law.
Check it exactly. If it fails, return any one connected original pair,
together with the free admissions. This remains a feasible answer.

On $\mathcal G$, for each selected path add a fresh source-copy leaf beyond
its original source occurrence leaf, and a fresh sink-copy leaf beyond its
original sink occurrence leaf. The two new edges have unit capacity. Extend
the selected path through them and give the resulting copy path flow
$\varepsilon$. These $K$ flows are jointly feasible: a normalized old edge
has load $\varepsilon Z_e\le1$, and a new edge has load $\varepsilon$.
Call the augmented graph $G^*$. The original occurrence edges remain present
with capacity one. All copies of original pair $i$ cross its same original
source occurrence edge.

\section{One decomposition and exact admission}\label{sec:transfer}

Give the $2K$ new copy terminals weight one and every other vertex weight
zero. Let $A$ be this integral weighting. Choose
\begin{equation}\label{eq:phi}
 \phi=\frac{\varepsilon}{8\Gamma(2K)}
      =\frac1{32T\Gamma(2K)}.
\end{equation}
Apply Proposition~\ref{prop:bjs} once. Its cut set satisfies
\begin{equation}\label{eq:cuts}
 |C|\le \phi(2K)\Gamma(2K)=\varepsilon K/4.
\end{equation}

\begin{lemma}\label{lem:charge}
At least $3K/4$ copy pairs have their entire prepared path in one component
of $G^*-C$.
\end{lemma}
\begin{proof}
Call a copy hit if its fixed prepared path uses a cut edge. The prepared
flows are feasible, so
\[
 \varepsilon\,\#\{\text{hit copies}\}
 \le\sum_{e\in C}\text{(prepared load on $e$)}
 \le |C|\le\varepsilon K/4.
\]
A path hit several times is only overcounted by the sum. An unhit path lies
entirely in one component. The charge is to individual path copies, without
discarding all copies of an original pair when just one is hit.
\end{proof}

For each component $U$, let $M_U$ be its unhit copy pairs and let $X_U$
be their endpoints. Ignore components with $M_U=\varnothing$. Work in the
induced graph $H_U=G^*[U]$, as promised by Proposition~\ref{prop:bjs}.
Although this can restore cut edges internal to $U$, the graphs $H_U$ have
disjoint edge sets for distinct components.

Write $q=|X_U|$. Demand $\phi/q$ on each unordered terminal pair has
marginal $\phi(q-1)/q\le\phi$. Dividing it by $\phi$ gives an
$A|_U$-respecting demand, because the endpoint indicator is pointwise at
most $A|_U$. Flow expansion and scaling therefore route the original
product demand with congestion one. A compact rational flow LP computes
such a routing. This argument directly uses all-demand expansion; it
does not assume that a normalized product-demand formula stays unchanged
when endpoints are deleted.

Proposition~\ref{prop:cks} applies to matching $M_U$ with $\alpha=\phi\le1$.
Summing its congestion-one routings over the edge-disjoint component graphs
selects at least
\begin{equation}\label{eq:yield}
 c_0\phi\sum_U|M_U|
 \ge \frac{3c_0K}{128T\Gamma(2K)}
 \ge \frac{3c_0F}{256\Gamma(2K)}
\end{equation}
unit copy demands. The real lower bound concerns an integer cardinality;
when it is smaller than a constant, one connected pair can be routed.

\begin{lemma}\label{lem:distinct}
Every congestion-one unit routing of copy pairs in $G^*$ selects at most
one copy of each indexed original pair. Projecting it to $G$ preserves
cardinality and feasibility.
\end{lemma}
\begin{proof}
For original pair $i$, its original source occurrence edge is a bridge of
the full augmented graph. Deleting it separates all source copies of $i$
from all sink copies of $i$, since the original endpoints are distinct.
Each admitted copy of $i$ sends one unit across this capacity-one cut.
Two admitted copies would force load at least two, which is impossible.
This cut argument applies to splittable flows. Indexed duplicates are
separate original pairs and have separate occurrence edges.

Remove cycles in the output flows and delete their copy and occurrence
segments. The resulting unit flows join their original endpoints and
use each original edge no more than before. Distinctness just proved
shows that cardinality is unchanged.
\end{proof}

\paragraph{The terminal parameter and running time.}
Since the present branch has $k>\lceil\log_2N\rceil$, we have
$T\le64(k+2)$. Equations~\eqref{eq:pieces} and $K\le D$ imply
$2K=O(k^2)$. Hence
\[
 \Gamma(2K)\le C'\log(k+2)
       \exp\!\bigl(C'\sqrt{\log\log(k+4)}\bigr)
\]
for an absolute $C'$. Combine this bound with~\eqref{eq:F},
\eqref{eq:yield}, and Lemma~\ref{lem:distinct}. The event $\mathcal G$ has
probability greater than $3/4$, and every branch returns feasible output.
If $z$ free pairs were removed initially, adding them preserves the
approximation: $z+r\ge(z+\OPT)/L$ whenever $r\ge\OPT/L$ and $L\ge1$.
The bound can be evaluated at the original number of pairs.

All augmentations, LPs, path decompositions, and fragment lists have
polynomial size and rational bit complexity. For completeness a suitable
integer overhead bound is
\[
 d(a)=\lceil\log_2(a+4)\rceil,\quad
 h(a)=\left\lceil\sqrt{\lceil\log_2(d(a)+2)\rceil}\right\rceil,\quad
 \Gamma(a)=C_0d(a)2^{C_1h(a)},
\]
for sufficiently large fixed positive integers $C_0,C_1$. It has the growth
in~\eqref{eq:gamma}, is nondecreasing, and is computed using integer bit
lengths and square roots. Thus $\phi$ is rational of polynomial bit length;
the constants are fixed constants of the cited decomposition, not advice
depending on the input. The exact small-$k$ branch uses at most $2N$ LPs;
the other branch uses the bounded sampler and the two polynomial-time
ingredients. There is no parameter-dependent exponent in the running time.
This completes the proof of Theorem~\ref{thm:main}.

\section{The exact-congestion rounding dependency}\label{sec:correction}

This section records how the strengthened clustering behind
Proposition~\ref{prop:cks} applies to our leaf endpoints. It explains an
established dependency, and does not assert a new nonuniform rounding theorem.
The relevant journal results are Lemmas~3.3--3.4 and~3.9--3.14
of~\cite{CKS13}, especially its Section~3.4, pp.~1478--1483.

Suppose a matching on $q$ leaf terminals has feasible product demand
$\alpha/q$. Give each terminal weight $\rho=\alpha/4$ and all other
vertices weight zero. For every ordered pair of terminals $u,v$, use
$\alpha/(4q)$ of the product flow in the corresponding direction; retain
the diagonal amount at $u$. Each terminal sends $\rho$ to the same uniform
distribution $\mu$ on the $q$ terminals. The two ordered directions use
half of each unordered commodity's given routing. The combined load is
therefore at most $1/2$ on each edge. This is the common-distribution
conclusion of journal Lemma~3.4 needed by the congestion-one rounding.
Its later clustering and tagging steps do not require the tree premise
used earlier to obtain that conclusion.

\paragraph{Coverage and ownership of clusters.}
Here is the residual-coverage convention for the leaf case. Root a spanning
tree at a zero-weight nonterminal. If total terminal weight is less than
one, route one matching pair, which already suffices for a constant-times-
$q\rho$ guarantee. Otherwise use the bottom-up grouping of journal
Lemma~3.3, with threshold one. It gives connected edge-disjoint clusters
with disjoint owned receiving sets of weight in $[1,2)$ and a residual
tree of unowned weight less than one. Absorb that residual tree into the
last extracted cluster, at its extraction root. Each earlier extracted
cluster meets the subsequently remaining tree in at most its extraction
root. Consequently there is full vertex coverage, induced clusters
remain edge-disjoint, and pairwise intersections have at most one vertex.
Owned masses are now in $[1,3)$. Positive-weight terminals are leaves,
so a shared connector has weight zero; receiving mass belongs only to
its owned set. This also ensures that every external tagging contact
belongs to a cluster.

Each cluster owns fewer than $3/\rho$ terminals. Greedily retaining a pair
and discarding pairs incident to either of its clusters discards at most
$6/\rho$ pairs per choice. Thus it keeps $\Omega(q\rho)$ pairs, and each
cluster contains at most the two terminals of one retained pair. Replacing
the usual upper mass two by three changes only an absolute constant.

\paragraph{Centers and tagging.}
We use the center, tagging, and pairing conclusions of journal
Lemmas~3.9 and~3.11--3.14. A center spreads one unit to its cluster's
receiving weights using at most half of each edge. A noncenter has a
separating cluster bridge, by the single-source min-cut characterization
in Lemma~3.12. Tag the first external contact along these separating cuts;
edges before this contact are bridges of the whole rounding-input graph.
Lemma~3.13 then makes the terminal a pseudocenter after joining an
appropriate auxiliary cluster: it spreads one unit within the receiving
weights using at most half capacity on nonbridges and at most one on
bridges of the whole input.

For clarity, the constant loss in this cited construction is independent
of $q$ and $\rho$. After identifying the two clusters of each chosen
pair, the tagging digraph has out-degree at most two on pair nodes.
The pair-node graph is 4-degenerate, so a five-coloring keeps a constant
fraction without tags to other retained pairs. Pairs whose own clusters
are already joined can be routed inside their disjoint amalgamations.
For the other pairs, transit clusters are unselected. Lemma~3.14 joins
tagging contacts through a transit spanning tree with edge-disjoint paths,
leaving at most one unmatched contact for odd degree at least three.
At least one third of pairs have neither endpoint isolated; a further
degree-two conflict selection keeps at least one third of those. These
are exactly the pseudocenter amalgamations and disjoint receiving supports
used by Lemma~3.11. A shared zero-weight vertex is a zero-edge tagging
contact, as specified in the last paragraph of journal Section~3.4.

\paragraph{Why the final congestion is one.}
Either a constant fraction of selected pairs have both endpoints in one
amalgamation, and have local edge-disjoint unit routes, or retain the
unmated class. In that class, spread each endpoint's unit to its owned
receiving support. Supports are disjoint, so the mass $\gamma(v)$ received
at terminal $v$ is at most $\rho$. Use $\gamma(v)/\rho$ times its
common-distribution flow constructed above. Their combined second-stage
load is at most $1/2$ per edge. Each endpoint reaches $\mu$; couple the
two distributions for each pair and concatenate the flows. Before cycle
removal, the loads are at most one on nonbridges and at most $3/2$ on bridges.

Journal Lemma~3.10 removes the apparent excess exactly. In a simple-path
decomposition, a bridge is used by a unit commodity with its whole unit
if the endpoints are separated, and with zero otherwise. Its total load
is an integer strictly below two, hence at most one. Cycle removal only
decreases loads elsewhere. All center tests, flow computations, grouping,
coloring, pairing, and cycle removal are polynomial. The common-distribution
premise, uniform owned masses, leaf connectors, and connected clusters
required by these cited steps have all been supplied. This is the
congestion-one dependency in Proposition~\ref{prop:cks} for our inputs.

\section*{Computational scope and AI methodology}

The algorithmic proof uses the cited polynomial-time decomposition and
rounding algorithms and rational multicommodity-flow LPs. Finite exact
rational checks accompanied development: they tested sampling marginals,
fixed-path cut charges, dyadic probability errors, copy-bridge separation,
and the small-parameter arithmetic on explicit fixtures. They are diagnostics,
not exhaustive graph tests or a replacement for the proof. The complete
Bansal--Jambulapati--Saranurak and Chekuri--Khanna--Shepherd algorithms have
not been implemented as a combined solver in this work.

OpenAI Codex was used substantially in proof development, adversarial
checking, primary-literature comparison, diagnostic Python code, and the
drafting and typesetting of all sections. The proof presentation and its
checks were materially generated or influenced by this assistance. Separate
AI-context checks do not provide formal proof certification. The named human
author remains responsible for the mathematical content and references.

\bibliographystyle{ACM-Reference-Format}
\bibliography{references}
\end{document}
